\section*{AI工具使用报告}
\subsection*{1. 对话可追溯性}
对于每次使用，提供准确的\textbf{提示词/查询}和\textbf{原始AI输出}（或具有代表性的摘录）。仅当AI仅用于翻译时才可省略完整日志；在这种情况下，需注明\textbf{工具名称和版本}。

\begin{itemize}[leftmargin=*]
    \item \textbf{示例---提示词：} "为模型部分建议一个大纲。" \textbf{输出（摘录）：} "1. 陈述假设；2. 定义符号；3. 推导主要方程\ldots"
    \item \textbf{仅使用翻译的团队：} 声明工具名称和版本（如DeepL 4.2、ChatGPT-4）；无需提供完整的提示词日志。
\end{itemize}

\subsection*{2. 分阶段声明}
说明AI在哪些阶段被使用，例如：
\begin{itemize}[leftmargin=*]
    \item \textbf{初始构思/头脑风暴}
    \item \textbf{代码修正/调试}
    \item \textbf{翻译/文字编辑}
\end{itemize}
标注每个阶段是否使用，并注明每个阶段使用的工具。

\subsection*{3. 风险验证}
描述团队如何检查\textbf{幻觉}和\textbf{错误引用}：例如，对关键声明进行人工事实核查、验证参考文献、交叉检验公式和代码。如果仅用于翻译，请声明没有将实质性技术内容委托给AI处理。
